% 前置部分：特别前言（1991）与原前言（1963）
% PDF p008（Special Preface）、p010–p011（Preface）

\chapter*{特别前言}
\addcontentsline{toc}{chapter}{特别前言}

我重读了《固体中的概念》（\emph{Concepts in Solids}），心情既是\emph{自豪}又是\emph{惭愧}。自豪，一方面因为正是这套讲义启发了 Brian Josephson 发明他以之命名的效应——不是每本书都能指明它所启发的那项确切的诺贝尔奖——另一方面也因为我的确在极短的篇幅内触及了某些关键论题，而这些论题至今在通常的固体理论教科书中仍未得到充分的讨论。比如说，令人震惊的是，这个国家使用的主流教材至今仍不触及 Mott 转变或"磁状态"（"Magnetic State"）。我瞄准的是\emph{概念性}的物理，而不是\emph{机械性}的物理，我希望我做到了。

惭愧，则是因为毕竟自那以来已经过去了 30 年的物理。例如我注意到，我断然摒弃紧束缚理论（tight-binding theory）实在是猜错了：它虽然尚未完全崭露头角，但以我目前的看法，它正是思考固体中大多数成键问题的正确途径。我并不为略去局域化（localization）而羞愧——当时只有 Mott 对它感兴趣，而我们俩当时都还不知道下一步该往哪里走。关于对称破缺（broken symmetry）我是有先见之明的——Josephson 也意识到了这一点——但我漏掉了相变，这一点我自己也指出过。

尽管如此，我相信一个普通学生读这本书所受的损害，仍会小于读随便哪几本我不便点名的其他书，因此我欢迎这次重印。

\begin{flushright}
P. W. Anderson\\
1991 年 3 月
\end{flushright}

\chapter*{前言}
\addcontentsline{toc}{chapter}{前言}

这些讲义是为 1961—1962 年秋冬学期在剑桥 Cavendish 实验室开设的一门课程写的。名义上，听课对象是修过固体物理概论、（至少）对理论感兴趣的二、三年级研究生；但我几乎不假定任何正式的\emph{理论}背景。我想，任何修过一门完整的量子力学课程的人都能读懂这些讲义；不过，对固体有所了解的读者会觉得容易得多，而且也不至于被我那些相当任意、相当专门的选材引入歧途。

这门课程的用意，是讲授固体物理的若干中心概念，并尽力挑选那些——能带理论、近自由电子、有效哈密顿量理论、元激发、对称破缺——在我看来离这门学科发展的主流尽可能近的概念。这样的选择必然是任意的：那些同样可以主张绝不次要、甚至同等重要的领域，如位错理论、输运理论与涨落-耗散定理、各种形式的磁共振理论，以及临界涨落，都被略去了，原因很简单——课程的长度是有限的。我对例子的选择就更加任意了。例如，选电场情形作为有效哈密顿量理论的例子、选激子作为集体激发的例子，是因为我想学生们在别处更可能遇到那些更常见的例子。为了给课程增加一点趣味，我不时引入一些\emph{原始材料}：关于近自由电子理论局限性的讨论、元激发理论的哲学的讨论，以及关于对称破缺的讨论，都是新的；而关于磁状态的那一部分则流传不广。

语言与讲述都非常不拘形式；与我当初写下的讲义原稿相比，改动极少。还要说明一句：我们几乎没有花力气把内容更新到当前。这两点局限，当然都隐含在"讲义卷"这个概念本身之中。

我愿感谢 Mott 教授与 Cavendish 实验室的盛情接待；实验室的秘书们整理了讲义的原始版本。Liu Sham 先生热心地为讲义做了编辑并补写了其中大部分公式。

\begin{flushright}
P. W. Anderson\\
Murray Hill, New Jersey\\
1963 年 1 月
\end{flushright}
